% TOP_PROBLEM: {"id": 3365, "title": "Are there an infinite number of lucky primes?", "category_id": 1, "category": {"id": 1, "name": "number_theory", "display_name": "Number Theory", "description": "Properties of integers, prime numbers, Diophantine equations.", "slug": "number-theory", "order_index": 1, "created_at": "2026-07-31T15:26:25.670Z"}, "status": "open", "problem_number": "OPG-37192", "set_id": 12, "statusQualification": "Repairs the corrupted sieve prose using the standard lucky-number steps 2,3,7,9,..., rather than consecutive step sizes."}
% FIRST_POSTED: 2026-10-01
% ENTRY_KIND: statement_only
% UnsolvedMath ID 3365; source catalogue code OPG-37192
% !TeX program = lualatex
% Definition and sources only; this file is not an LLM solution attempt.
\documentclass[11pt,a4paper]{article}
\usepackage[margin=25mm]{geometry}
\usepackage{fontspec}
\setmainfont{lmroman10-regular.otf}[BoldFont=lmroman10-bold.otf,
 ItalicFont=lmroman10-italic.otf,BoldItalicFont=lmroman10-bolditalic.otf]
\usepackage{amsmath,amssymb,mathtools}
\usepackage{unicode-math}
\setmathfont{latinmodern-math.otf}
\AtBeginDocument{\let\setminus\smallsetminus}
\usepackage{xurl,hyperref}
\hypersetup{colorlinks=true,urlcolor=blue}
\setcounter{secnumdepth}{0}
\setlength{\parindent}{0pt}
\setlength{\parskip}{5pt}
\setlength{\emergencystretch}{3em}
\begin{document}
\section*{Are there an infinite number of lucky primes?}\label{3365}
{\small UnsolvedMath ID 3365 · OPG-37192 · Number Theory · OpenGarden}
\subsection{Definitions and mathematical statement}
Begin with the increasing sequence $L_1=(1,3,5,7,9,\ldots)$ obtained by deleting every second entry from the positive integers. For each integer $j\ge2$, let $\ell_j$ be the $j$th entry of the increasing sequence $L_{j-1}$. Form $L_j$ by deleting entries whose positions in $L_{j-1}$ are multiples of $\ell_j$, with positions counted from one and all deletions in that stage determined simultaneously. Also put $\ell_1=1$.

The lucky numbers are the positive integers surviving every stage, equivalently the increasing sequence $\ell_1,\ell_2,\ldots$. Its beginning is
\[
                         1,3,7,9,13,15,21,25,31,\ldots.
\]
After the initial deletion of even positions, the successive sieve step sizes are therefore $3,7,9,\ldots$, chosen from the surviving sequence itself. They are not the consecutive integers $3,4,5,\ldots$ and are not necessarily prime.

A lucky prime is a lucky number that is prime in the ordinary multiplicative sense: it is greater than one and has no positive divisors other than one and itself. The conjecture asks whether there are infinitely many lucky primes. Equivalently, for every positive integer $B$ there should exist a surviving number $\ell>B$ with this primality property.

Infinitely many lucky numbers alone do not establish the conjecture, since surviving the positional sieve is not a divisibility test. The source's prose mixes these sieve conventions; the recursive definition above fixes the standard lucky-number sequence intended by its title and discussion.
\subsection{Short English statement}
Are infinitely many numbers surviving the standard lucky-number positional sieve ordinary prime numbers?
\subsection{Sources}
\begin{itemize}
\item[S1] ULAM.AI and the UnsolvedMath Contributors, supplied problems.json, record 3365 (OPG-37192), OpenGarden. Catalogue identity and source transcription; CC BY 4.0.
\url{https://huggingface.co/datasets/ulamai/UnsolvedMath}.
\item[S2] Open Problem Garden, Are there an infinite number of lucky primes?. Original problem and discussion.
\url{http://www.openproblemgarden.org/op/are_there_an_infinite_number_of_lucky_primes}.
\item[S3] C. Sanna, Explicit inequalities for the nth lucky number (2026).
\url{https://arxiv.org/abs/2604.07142}.
\end{itemize}
\end{document}
